\begin{resumo}[\protect\bfseries Resumo]
  No presente documento é apresentada a experiência do aluno Guilherme Argilar durante sua passagem pela Agência Experimental de Engenharia de Software (AGES) da PUCRS, ambiente essencial para o intercâmbio de conhecimento e o aprendizado prático no curso de Bacharelado em Engenharia de Software. A AGES é composta por quatro disciplinas oferecidas ao longo da graduação, nas quais os alunos assumem responsabilidades progressivamente maiores conforme sua senioridade: como AGES I, executam tarefas de programação; como AGES II, atuam na análise de dados da aplicação; como AGES III, no desenvolvimento da arquitetura de software; e, como AGES IV, na gestão da equipe do projeto. Ao longo de sua graduação, o autor participou dos quatro módulos da Agência, vivenciando diferentes papéis no processo de desenvolvimento de software e na dinâmica do trabalho em equipe. Este documento narra sua jornada ao longo do curso, com foco nas experiências dos semestres 2023/1, 2024/1, 2025/1 e 2026/1, durante os quais esteve envolvido nos projetos Focus, Decola, Totem Interativo e Estudify, respectivamente. Em cada módulo, são destacados os aprendizados, os desafios enfrentados e a evolução das responsabilidades dentro da AGES.
  \bigbreak\
  \\\textbf{PALAVRAS CHAVES:}
  AGES, Engenharia de Software, Aprendizado, Processo de Desenvolvimento de Software, Projeto.
\end{resumo}